% from ega1-6-fr.tex
\oldpage[I]{155}
est contenu dans $f^{-1}(V)$ et de type fini sur $V$~; par suite on a
$g(f(U))\subset W$ et $U$ est de type fini sur $W$
(\hyperref[I.6.3.4-fr]{6.3.4, (ii)}). Enfin, pour démontrer (iii), on peut se
borner au cas où $Y=S$ (\hyperref[I.3.3.11-fr]{3.3.11})~; pour tout
$x'\in X'=X_{(S')}$, soient $x$ l'image de $x'$ dans $X$, $s$ l'image de $x$
dans $S$, $T$ un voisinage ouvert de $s$, $T'$ son image réciproque dans $S'$,
$U$ un voisinage ouvert de $x$ dont l'image est contenue dans $T$ et qui est de
type fini sur $T$~; alors $U\times_S T'=U\times_T T'$ est un voisinage ouvert
de $x'$ (\hyperref[I.3.2.7-fr]{3.2.7}) qui est de type fini sur $T'$
(\hyperref[I.6.3.4-fr]{6.3.4, (iv)}).

\begin{corollary}[6.6.7]
\label{I.6.6.7-fr}
Soient $X$, $Y$ deux $S$-préschémas qui sont localement de type fini sur $S$.
Si $S$ est localement noethérien, $X\times_S Y$ est localement noethérien.
\end{corollary}

En effet, $X$ étant localement de type fini sur $S$, est localement
noethérien, et $X\times_S Y$ est localement de type fini sur $X$, donc est
aussi localement noethérien.

\begin{remark}[6.6.8]
\label{I.6.6.8-fr}
La proposition (\hyperref[I.6.3.10-fr]{6.3.10}) et sa démonstration s'étendent
aussitôt au cas où l'on suppose seulement que le morphisme $f$ est localement
de type fini. De même, les propositions
(\hyperref[I.6.4.2-fr]{6.4.2}) et (\hyperref[I.6.4.9-fr]{6.4.9}) restent
valables lorsqu'on suppose que les préschémas $X$, $Y$ qui figurent dans leur
énoncé sont seulement localement de type fini sur le corps $K$.
\end{remark}

% (continued in ega1-7-fr.tex, before its first page marker)
\section{Applications rationnelles}
\label{section:I.7-fr}

\subsection{Applications rationnelles et fonctions rationnelles.}
\label{subsection:I.7.1-fr}

\begin{env}[7.1.1]
\label{I.7.1.1-fr}
Soient $X$, $Y$ deux préschémas, $U$, $V$, deux ouverts denses dans $X$, $f$
(resp. $g$) un morphisme de $U$ (resp. $V$) dans $Y$~; nous dirons que $f$ et
$g$ sont
\emph{équivalents} s'ils coïncident dans un ouvert dense dans $U\cap V$.
Comme une intersection finie d'ouverts denses dans $X$ est un ouvert dense
dans $X$, il est clair que cette relation est une
\emph{relation d'équivalence}.
\end{env}

\begin{definition}[7.1.2]
\label{I.7.1.2-fr}
Étant donnés deux préschémas $X$, $Y$, on appelle
\emph{application rationnelle de $X$ dans $Y$} une classe d'équivalence de
morphismes de parties ouvertes denses de $X$ dans $Y$, pour la relation
définie dans (\hyperref[I.7.1.1-fr]{7.1.1}). Si $X$ et $Y$ sont des
$S$-préschémas, on dit qu'une telle classe est une
\emph{$S$-application rationnelle} s'il existe un représentant de cette classe
qui est un $S$-morphisme. On appelle
\emph{$S$-section rationnelle de $X$} toute $S$-application rationnelle de $S$
dans $X$. On appelle
\emph{fonction rationnelle sur un préschéma $X$} toute $X$-section rationnelle
du $X$-préschéma $X\otimes_{\mathbf{Z}}\mathbf{Z}[T]$ (où $T$ est une
indéterminée).

Par abus de langage, lorsqu'il ne sera question que de $S$-préschémas, on dira
«~application rationnelle~» au lieu de «~$S$-application rationnelle~» si
aucune confusion ne peut en résulter.

Soient $f$ une application rationnelle de $X$ dans $Y$, $U$ un ouvert de
$X$~; si $f_1$, $f_2$ sont des morphismes appartenant à la classe $f$, définis
respectivement dans des ouverts denses $V$, $W$ de $X$, les restrictions
$f_1|(U\cap V)$, $f_2|(U\cap W)$ coïncident dans $U\cap V\cap W$ qui est dense
dans $U$~; la classe de morphismes $f$ définit donc une application rationnelle
de $U$ dans $Y$, appelée
\emph{restriction de $f$ à $U$} et notée $f|U$.

Si, à tout $S$-morphisme $f:X\to Y$, on fait correspondre la $S$-application
rationnelle
